
\chapter{Grover against the chain}

\textbf{Purpose:} Price the one quantum algorithm that genuinely applies to unstructured search,
against the live difficulty, and not in the abstract, and separate the reduction it earns from
the advantage it does not. \textbf{Scope:} \texttt{qubit\_trace.py},
\texttt{nonce\_period.py --grover}, \texttt{nonce\_period.py --dimensional}.

Grover is the correct algorithm for unstructured search and it gives a genuine square root. The
algorithm was already implemented here before the question was asked of it; what was missing was the
accounting.

\section{The implementation, and that it works}

\texttt{qubit\_trace.py} runs Grover with amplitude amplification and an optimal-time detector.
\(P(\text{marked})\) climbs from \(0.003906\) to \(0.999947\) at step 12, is optimal at 13, and then
\textbf{falls back}. The fall-back is the clearest evidence available that the iteration
rotates, and does not accumulate: an accumulating process does not come back.
\((\pi/4)\sqrt{256} = 12.57\), and the optimal count is therefore right.

\section{The reduction is real, and it is huge}

Against block 966463 at difficulty \(1.2745 \times 10^{14}\):

\begin{center}
\begin{tabular}{lr}
\toprule
classical hashes for one block & \(5.474 \times 10^{23}\) \\
Grover oracle calls, \((\pi/4)\sqrt{N}\) & \(5.811 \times 10^{11}\) \\
\textbf{reduction in calls} & \(\mathbf{9.420 \times 10^{11}}\) \\
\bottomrule
\end{tabular}
\end{center}

\section{What a Grover machine is worth in classical hash rate}

This is the only comparison here that is apples to apples, and the conversion is exact. A machine at
\(G\) oracle calls a second finds a block in \(\text{iterations}/G\) seconds; a classical miner at
\(H\) hashes a second finds one in \(N/H\) seconds. So the machine equals
\(G \times 9.420 \times 10^{11}\) hashes a second.

\begin{center}
\begin{tabular}{rrrrr}
\toprule
oracle calls a s & equivalent H/s & vs this card & vs network & time to a block \\
\midrule
1 & \(9.420 \times 10^{11}\) & \(\mathbf{399\times}\) & \(1.03 \times 10^{-9}\) & \(1.84 \times 10^{4}\) years \\
\(10^{4}\) & \(9.420 \times 10^{15}\) & \(3.99 \times 10^{6}\) & \(1.03 \times 10^{-5}\) & 1.84 years \\
\(10^{7}\) & \(9.420 \times 10^{18}\) & \(3.99 \times 10^{9}\) & \(1.03 \times 10^{-2}\) & 0.7 days \\
\bottomrule
\end{tabular}
\end{center}

\textbf{Grover is not useless and an earlier draft of this work implied it was.} At even one oracle
call a second the machine is worth \(9.420 \times 10^{11}\) hashes a second, about four hundred times
the card in this room. The square root does real work and the algorithm choice was correct.

What it loses to is the \textbf{network} and not the card. At \(9.123 \times 10^{20}\) hashes a
second the network outruns a one-call-a-second machine by \(9.68 \times 10^{8}\).

\section{Why the oracle is the obstruction}

A Grover call is not a hash. It is a coherent reversible circuit for the whole function, run in
superposition, with no measurement until the end. \emph{Quoted, and not verified here:} published
estimates put a quantum SHA-256 circuit at a few thousand logical qubits and \(10^{5}\) to
\(10^{6}\) T gates per evaluation, and mining needs SHA-256 twice.

\begin{center}
\begin{tabular}{lrr}
\toprule
logical T gate & T per oracle & oracle calls a second \\
\midrule
1 \(\mu\)s & \(10^{5}\) & \(1.000 \times 10^{1}\) \\
1 ns & \(10^{5}\) & \(1.000 \times 10^{4}\) \\
1 ps & \(10^{5}\) & \(1.000 \times 10^{7}\) \\
\bottomrule
\end{tabular}
\end{center}

Against the \(9.685 \times 10^{8}\) needed for a ten-minute block. A picosecond logical T gate is far
beyond anything projected and still falls two orders short.

\section{The decisive point needs no hardware estimate}

Mining is a race between many machines, and the two types scale oppositely when more are added:

\begin{center}
\begin{tabular}{rrr}
\toprule
machines & classical & Grover \\
\midrule
10 & \(10\times\) & \(3.2\times\) \\
10\,000 & \(10\,000\times\) & \(100\times\) \\
1\,000\,000 & \(1\,000\,000\times\) & \(1\,000\times\) \\
\bottomrule
\end{tabular}
\end{center}

Parallel Grover splits the search space, and \(m\) machines therefore buy \(\sqrt{m}\). Classical search buys
\(m\). \textbf{The square root that helps a single searcher hurts a fleet}, and this argument
survives any hardware estimate because it contains none.

\section{What partial information would be worth}

If the space factored into \(d\) independently searchable dimensions, Grover would cost
\(d\,N^{1/2d}\):

\begin{center}
\begin{tabular}{rrr}
\toprule
dimensions & oracle calls & vs plain Grover \\
\midrule
1 & \(7.399 \times 10^{11}\) & \\
8 & \(2.436 \times 10^{2}\) & \(2.39 \times 10^{9}\) \\
\textbf{32} & \(\mathbf{75}\) & \(\mathbf{7.73 \times 10^{9}}\) \\
\bottomrule
\end{tabular}
\end{center}

\textbf{Seventy-five oracle calls.} The arithmetic is correct. It needs partial credit: an
oracle reporting \emph{this dimension is right} while the others are wrong. A hash function is built
to destroy exactly that, since every input bit reaches every output bit.

Four instruments here looked for partial credit independently and none found it:
\texttt{nonce\_read.py} null with a drawn null over 200 shuffles; \texttt{beam\_read.py} null at
every live round with 256 line integrals and a weakest limit of 0.100; \texttt{--invert}, where a
partial inversion returns zero at best and a loss at worst; and \texttt{--nested}, which shows no
partial credit on the way and reach growing as \(b^d\).

\emph{Quoted, and not verified here:} the BBBV bound proves \(\sqrt{N}\) is optimal for a black-box
oracle, any scheme beating it exploits structure. This one beats it by \(7.73 \times 10^{9}\),
the measure of how much structure it silently assumes.

\section{Nest and invert is the price of coherence}

Computing forward, marking, then running the circuit backwards to uncompute is Bennett's
construction and it is \textbf{mandatory}: a Grover oracle must leave its ancillas clean or the
interference does not happen. An oracle therefore costs roughly twice the forward circuit.

Measured here, and not argued. \texttt{--invert} returns local purity to \(1.000000000000\) from
\(0.890625\) when the inversion reaches the whole structure, and a partial inversion returns zero at
best and a loss at worst. \texttt{--nested} shows reach growing as \(b^d\) with no partial credit on
the way. Full reach works and partial reach loses: the same sentence as the section above.
